\section{Introduction}

% ------------------------------------------------------------
\begin{frame}
  \frametitle{About the Lecturer}

  \begin{block}{Mathis Neunzig}
    \begin{itemize}
      \item B.Sc. Business Information Technology (DHBW Mannheim)
      \item M.Sc. Data Science (Fernuniversität Hagen)
      \item ABAP Developer \& Trainer at SAP
      \item Teaches Databases \& Software Engineering at DHBW
      \item Based in Mannheim (Baden-Württemberg)
      \item Originally from Oldenburg (Lower Saxony)
    \end{itemize}
  \end{block}

  \begin{block}{Contact}
    \footnotesize
    \textbf{Email}\\
    mathis.neunzig@gmail.com\\[6pt]
    \textbf{Phone}\\
    +49\,174\,9885992
  \end{block}
\end{frame}

% ------------------------------------------------------------
\begin{frame}
  \frametitle{Getting to Know You}

  \begin{block}{Let's go around the room -- please tell us:}
    \begin{enumerate}
      \item Your \textbf{name}
      \item Your \textbf{employer}
      \item Your \textbf{favourite animal}
      \item Your \textbf{favourite song, album or artist} right now
      \item Your \textbf{favourite programming language} (so far)
    \end{enumerate}
  \end{block}
\end{frame}

% ------------------------------------------------------------
\begin{frame}
  \frametitle{Agenda}

  This course is structured around multiple interconnected topic groups that build on each other progressively throughout the semester.

  \begin{block}{1. Foundations}
    \begin{itemize}
      \item Databases and Database Management Systems (DBMS)
      \item Data models
      \item Relational algebra
      \item Relational databases
      \item Data modelling \& the DB design process
    \end{itemize}
  \end{block}

  \begin{block}{2. SQL}
    \begin{itemize}
      \item Data Manipulation Language (DML)
      \item Data Definition Language (DDL)
      \item Data Control Language (DCL)
    \end{itemize}
  \end{block}
\end{frame}

\begin{frame}
  \begin{block}{3. Advanced Topics}
    \begin{itemize}
      \item Database transaction management
      \item Concurrency control
      \item Performance optimisation
      \item Backup \& Recovery
      \item Security \& SQL Injection
      \item Object-Relational Mapping (ORM)
    \end{itemize}
  \end{block}

  \begin{block}{4. DB Project (4th semester)}
    \begin{itemize}
      \item Recap of DB design process
      \item Recap of ORM
      \item Case study: designing a database for a real-world scenario
    \end{itemize}
  \end{block}
\end{frame}

\begin{frame}
  \frametitle{Semester Schedule}

  \scriptsize
  \begin{tabularx}{\linewidth}{@{}l l l X@{}}
    \toprule
    \textbf{Session} & \textbf{Date} & \textbf{Time} & \textbf{Topics} \\
    \midrule
    1  & 09.09 & 12:30--15:00 & Introduction; Course overview; Exam info; DB Management System; Data Models \\
    2  & 15.09 & 14:00--17:15 & Relational Databases; Basic Relational Algebra \\
    3  & 19.09 & 09:00--12:15 & Advanced Relational Algebra; Normal Forms (1NF, 2NF, 3NF) \\
    4  & 26.09 & 09:00--12:15 & Normal Forms (BCNF, 4NF, 5NF, 6NF) \\
    5  & 29.09 & 14:00--17:15 & ERM Design \\
    6  & 13.10 & 14:00--17:15 & SQL (DML) \\
    7  & 20.10 & 14:00--17:15 & SQL (DML) \\
    8  & 27.10 & 14:00--17:15 & SQL (DCL); CAP-Theorem \\
    9  & 03.11 & 14:00--17:15 & Backup \& Recovery; Performance; Concurrency \\
    10 & 10.11 & 14:00--17:15 & Security; Object-Relational Mapping \\
    11 & 17.11 & 14:00--17:15 & Exam preparation \\
    12 & 26.11 & 10:00--11:00 & \textbf{Written Exam} \\
    \bottomrule
  \end{tabularx}

\end{frame}

% ------------------------------------------------------------
\begin{frame}
  \frametitle{Organization}

  \begin{block}{Mailing List \& Contact}
    \begin{itemize}
      \item All course announcements will be sent via discord.
      \item There is no mailing list in place for this class. Please check the discord channel regularly for updates.
      \item Lecturer contact: \texttt{mathis.neunzig@gmail.com} \quad +49\,174\,9885992
    \end{itemize}
  \end{block}

  \begin{alertblock}{Emergency Contact}
    If you have an urgent issue that cannot wait for the next session -- for example a technical problem on the day of an assignment submission -- please use the phone number above. For non-urgent matters, emailor discord is always preferred.
  \end{alertblock}
\end{frame}

% ------------------------------------------------------------
\begin{frame}
  \frametitle{Rules}

  \begin{block}{General Conduct}
    Standard DHBW rules of conduct apply throughout all sessions.
  \end{block}

  \begin{itemize}
    \item \textbf{Absences:} Please give advance notice if you cannot attend. A quick message the day before is always appreciated. If you miss a session, it is your responsibility to catch up on the material.
    \item \textbf{Deadlines:} If you foresee that you cannot meet an assignment deadline, contact the lecturer \emph{before} the deadline, not after. Extensions can often be arranged if asked in advance.
    \item \textbf{Open door:} Feel free to reach out at any time with questions, concerns, or feedback; about the content, the pace, or anything else. There are no silly questions in this course.
    \item \textbf{Group exercises:} Active participation during group exercises is expected and contributes to your own learning as well as the group's. Sitting passively during a group task is not acceptable.
  \end{itemize}

\end{frame}

% ------------------------------------------------------------
\begin{frame}
  \frametitle{Assessment: Written Exam}

  \begin{block}{Key Facts}
    \begin{itemize}
      \item \textbf{Weight:} 40 points \quad(40\,\% of final grade)
      \item \textbf{Timing:} End of 3rd semester (12, 26.11)
      \item \textbf{Duration:} 60 minutes
    \end{itemize}
  \end{block}

  \begin{alertblock}{Exam Scope}
    Topics marked with \textbf{(!)} in the slides will \emph{not} appear in the written exam. These are advanced or supplementary topics included for completeness and professional context, but you will not be tested on them under exam conditions.
  \end{alertblock}

  \begin{examples}{Permitted Aids}
    You may bring \textbf{one handwritten A4 sheet} (both sides allowed) as a personal cheat sheet. It must be handwritten -- printed or typed sheets are not permitted. Writing your own cheat sheet is itself an effective revision strategy.
  \end{examples}

\end{frame}

% ------------------------------------------------------------
\begin{frame}
  \frametitle{Assessment: Assignments and Project}

  \begin{block}{Overview -- 60 points total (60\,\% of final grade)}
    \begin{tabularx}{\linewidth}{@{}l X r@{}}
      \toprule
      \textbf{Assignment} & \textbf{Description} & \textbf{Points} \\
      \midrule
      \multicolumn{3}{@{}l}{\textit{Semester 3}} \\
      Assignment 1 & Individual work & 10 pts \\
      Assignment 2 & Individual work & 10 pts \\
      \multicolumn{3}{@{}l}{\textit{Semester 4}} \\
      Project & Final project (in groups) & 40 pts \\
      \bottomrule
    \end{tabularx}
  \end{block}

  \begin{block}{Submission Policy (Assignments 1 and 2)}
    \begin{itemize}
      \item \textbf{1st submission} receives a grade \emph{and} written feedback.
      \item \textbf{Optional 2nd submission} is possible after reviewing the feedback.
      \item The 2nd submission grade replaces the 1st -- use this opportunity wisely.
    \end{itemize}
  \end{block}

\end{frame}

% ------------------------------------------------------------
\begin{frame}
  \frametitle{Literature}

  \begin{alertblock}{No textbook purchase is required to pass this course.}
    All essential material is covered in the lecture slides and exercises.
  \end{alertblock}

  \textbf{Recommended reading (for those who want to go deeper):}

  \begin{itemize}
    \item Kemper \& Eickler -- \textit{Datenbanksysteme}\\
          \footnotesize The German-language standard reference for database courses at German universities. Thorough and mathematically precise.\normalsize

    \item Garcia-Molina, Ullman \& Widom -- \textit{Database Systems: The Complete Book}, 2nd~ed.\\
          \footnotesize Comprehensive English-language textbook covering theory and practice in depth.\normalsize

    \item Foster \& Godbole -- \textit{Database Systems: A Pragmatic Approach}, 2nd~ed.\\
          \footnotesize A more application-oriented introduction, well-suited as a companion to this course.\normalsize

    \item Gumm \& Sommer -- \textit{Einführung in die Informatik}, 10th~ed.\\
          \footnotesize Broader CS introduction, useful for contextualising databases within the wider field.\normalsize
  \end{itemize}

\end{frame}
